\scanpage{53}
\mspage{43}
\margin{n°382}

4.13 -- \uline{Affaiblissement de structure}.

Dans tout ce numéro, on suppose que $K=\uuline{R}$ ; on note I l'ensemble $\{1,2,\ldots,n,\ldots;\infty;\omega\}$. Nous ordonnerons I par la relation d'ordre usuelle sur les entiers, étendue à $\infty$ et à $\omega$ au moyen des formules
\[
n<\infty<\omega.
\]

4.13.1. Soient $r,s\in I$, avec $r\leq s$, et soit X une variété de classe $C^s$. Il existe sur l'espace topologique X une structure de variété de classe $C^r$ et une seule telle que toute carte de X pour la structure donnée soit une carte de X pour cette nouvelle structure. Soit $X_r$ la variété de classe $C^r$ ainsi définie. On dit qu'elle se déduit de X par \uline{affaiblissement} de la structure de variété de X ; ou encore que sa structure de variété est \uline{sous-jacente} à celle de X. La notion d'affaiblissement est transitive : si $r'\leq r$, on a $X_{r'}=(X_r)_{r'}$ ; elle commute aux produits : si Y est de classe $C^s$, on a $(X\times Y)_r=X_r\times Y_r$ ; on a un résultat analogue pour $X\times_S Y$ sous les hypothèses du n° 4.11.3.

Soit $a\in X$ et soit c une carte de X en a ; c'est aussi une carte de $X_r$ en a. On en déduit (cf. n° 4.5.1) des isomorphismes
\[
\vartheta_c:K^n\longrightarrow T_a(X),\qquad
\vartheta'_c:K^n\longrightarrow T_a(X_r),\qquad n=\dim_x(X),
\]
d'où un isomorphisme $\vartheta'_c\circ\vartheta_c^{-1}:T_a(X)\longrightarrow T_a(X_r)$. Cet isomorphisme est indépendant du choix de c ; il permet \uline{d'identifier les espaces tangents} à X et à $X_r$ en a.

4.13.2. Soit X (resp. X') une variété de classe $C^s$ (resp. $C^{s'}$), et soit $r\in I$, $r\leq\operatorname{Inf}(s,s')$. Une application $f:X\longrightarrow X'$ est dite \uline{de classe} $C^r$ si c'est un morphisme de $X_r$ dans $X'_r$ ; une telle application est de classe $C^{r'}$ pour tout $r'\leq r$. De plus, l'application linéaire tangente à $f:X_r\longrightarrow X'_r$ en un point $a\in X$ coïncide avec l'application linéaire tangente à f considérée comme morphisme de $X_{r'}$ dans $X'_{r'}$. Pour que $f:X_r\longrightarrow X'_r$ soit une submersion (resp. une immersion, une subimmersion), il faut et il suffit que $f:X_{r'}\longrightarrow X'_{r'}$ en soit une ; on dira alors plus simplement que f est une submersion (resp. une immersion, resp. une
\mspage{52}
